\chapter{Implementación}\label{cap:implementacion}

\section{Estructura del proyecto y tecnologías utilizadas}

La implementación de GlucoMate se divide en dos grandes subsistemas: un
\textbf{frontend móvil} desarrollado con React Native y Expo, y un \textbf{backend}
basado en Node.js, Express y PostgreSQL. A partir del proyecto adjunto, la estructura del
repositorio puede resumirse del siguiente modo:

\begin{itemize}
    \item \texttt{frontend/}: aplicación móvil Android, navegación, pantalla principal,
          perfil y modales de interacción.
    \item \texttt{backend/}: API REST, acceso a base de datos y lógica de integración con
          servicios externos.
\end{itemize}

En el frontend se emplean React Native, Expo, \texttt{react-native-nfc-manager},
\texttt{react-native-svg}, D3.js y \texttt{@react-navigation/bottom-tabs}. En el backend
se utilizan Node.js, Express, \texttt{pg} para PostgreSQL y el SDK de Groq para
la integración con el modelo generativo \texttt{llama-3.3-70b-versatile}. Esta selección tecnológica refleja una decisión
de diseño clara: mantener una interfaz móvil moderna y, al mismo tiempo, desacoplar del
cliente toda la lógica sensible y la persistencia estructurada.

\section{Módulo NFC --- \texttt{NFCScanner.js}}

El fichero \texttt{NFCScanner.js} concentra la lógica principal de la aplicación móvil.
Desde este componente se inicia el proceso de lectura del sensor, se solicitan datos al
puente de comunicación utilizado y se actualiza tanto la gráfica histórica como el valor
actual de glucosa. También coordina la apertura de modales y la recarga del historial.

La implementación real confirma este papel central. El componente inicializa el gestor
NFC al montarse, comprueba si la tecnología está soportada y cancela cualquier petición
pendiente al desmontarse. Además, mantiene en memoria el histórico glucémico, el perfil,
los eventos clínicos recientes y el estado de visibilidad de los distintos modales.

El proceso de escaneo sigue una secuencia concreta: solicitud de tecnología
\texttt{NfcTech.NfcV}, cancelación controlada de la petición, recuperación del dato desde
Juggluco y posterior envío de la lectura al backend mediante \texttt{POST} a
\texttt{/api/glucose}. Tras guardar la lectura, el frontend solicita de nuevo el histórico
de 12 horas para refrescar la visualización.

\section{Módulo de visualización}

La visualización integra el valor actual de glucosa con una representación histórica de
las últimas horas. Para ello se apoya en D3.js y \texttt{react-native-svg}, lo que
permite dibujar curvas, líneas de referencia y marcadores de eventos clínicos dentro de
la misma pantalla.

En la implementación observada, el gráfico se construye a partir del conjunto devuelto por
\texttt{/api/glucose/history}. Ese conjunto se transforma en una estructura temporal apta
para D3, sobre la que se calculan escalas, trayectorias y posiciones de cada evento. La
pantalla superpone además múltiples capas de información: curva de glucosa, indicador del
valor actual, tendencia, insulina activa, comidas registradas y actividad física reciente.

La decisión de usar SVG nativo resulta especialmente adecuada en este contexto, ya que
permite controlar con precisión cada elemento visual y combinar datos continuos
(trayectorias) con datos discretos (eventos y marcadores) dentro del mismo lienzo.

\section{Módulo de registro clínico}

Los modales \texttt{InsulinModal.js}, \texttt{FoodModal.js} y
\texttt{ExerciseModal.js} encapsulan el registro de insulina, comidas y actividad
deportiva. Cada uno recoge los datos necesarios, los valida y los envía al backend para
su persistencia.

El backend expone para ello los endpoints \texttt{/api/insulin-events},
\texttt{/api/food-events} y \texttt{/api/exercise-events}, tanto en lectura como en alta.
En particular, el flujo de comidas presenta un interés adicional porque no solo registra
hidratos de carbono, sino que puede acoplar el registro de un bolo asociado dentro de la
misma operación lógica. Esto refuerza la idea de que la aplicación no almacena eventos
aislados, sino relaciones clínicamente significativas entre ellos.

\section{Módulo de perfil de usuario}

La pantalla \texttt{ProfileScreen.js} permite almacenar la configuración terapéutica del
usuario, incluyendo insulinas, ratios ICR por franja horaria y factor de sensibilidad.
Estos datos constituyen la base del comportamiento personalizado del sistema.

Desde el punto de vista de implementación, el perfil se apoya en dos endpoints:
\texttt{GET /api/profile} para recuperar la configuración actual y
\texttt{PUT /api/profile} para actualizarla. La interfaz también consulta
\texttt{/api/insulins} para poblar los selectores de insulina rápida y basal. El resultado
es un formulario relativamente compacto, pero con un impacto importante sobre el resto del
sistema, ya que esos parámetros son reutilizados por el backend y por el agente de IA.

\section{Agente de IA --- lógica de recomendación}

La funcionalidad de apoyo inteligente se articula sobre la información clínica del
usuario y su contexto reciente. La lógica de recomendación utiliza el perfil terapéutico,
la glucosa actual, la insulina activa y el historial inmediato para producir respuestas
orientativas, siempre bajo un enfoque de apoyo y no de sustitución del criterio médico.

La implementación del backend muestra una estrategia razonable de construcción de
contexto. Antes de invocar el modelo, el servidor consulta en paralelo el perfil del
usuario, la insulina activa de las últimas horas, el ejercicio reciente y la última
lectura de glucosa. Posteriormente calcula reducciones porcentuales en función del tipo de
deporte, detecta situaciones de advertencia y compone un contexto textual que se remite al
modelo generativo. Esta decisión traslada al servidor la parte más sensible de la lógica y
evita que el cliente deba contener reglas clínicas o credenciales externas.

\section{Integración y flujo global de la app}

El flujo completo de implementación puede resumirse así: lectura del sensor, recepción de
datos, persistencia en backend, recuperación del histórico, renderizado en interfaz y,
cuando procede, generación de recomendaciones contextualizadas.

En términos de integración, la aplicación presenta una arquitectura en cadena: sensor,
aplicación puente, frontend móvil, backend y base de datos. No es una arquitectura
puramente local, pero tampoco depende todavía de un despliegue cloud completo. Esto encaja
con el carácter de prototipo funcional del TFG: suficientemente realista para validar el
flujo extremo a extremo, pero todavía acotado para mantener la complejidad bajo control.